\section{回顾与动机}

本讲是 DDPM（Denoising Diffusion Probabilistic Models）系列的第二部分。在上一讲中，我们了解了扩散模型的基本框架：从 VAE 的局限性出发，引入多层潜变量的思路，进而定义了扩散模型的前向过程和反向过程。本讲将从训练目标的推导开始，逐步推导出简化的训练损失函数，并引入三种等价的参数化方式：均值预测、干净样本预测和噪声预测。

\subsection{从 VAE 到扩散模型的动机}

回顾生成模型的核心目标：建模边际数据分布 $p(x_0)$。在 VAE 框架中，这通过先验分布 $p(z)$（标准高斯）和似然分布 $p_\theta(x|z)$（解码器）的组合来实现：

$$p_\theta(x_0) = \int p(z) \, p_\theta(x_0 | z) \, dz$$

VAE 的局限在于，先验和似然都被建模为高斯分布，这限制了边际分布 $p_\theta(x_0)$ 的表达能力——它本质上是高斯分布的混合，可能无法充分捕捉复杂数据分布的多模态特性。

\begin{importantbox}{扩散模型的核心思想}
扩散模型通过引入\textbf{多步}潜变量序列 $x_1, x_2, \dots, x_T$ 来增强似然分布的表达能力。虽然每一步的转移分布仍然是高斯分布，但通过多步非线性变换的组合，整体的似然分布可以表达远比单个高斯更复杂的分布。
\end{importantbox}

\subsection{扩散模型的基本设定}

扩散模型的核心组件包括：

\begin{itemize}
    \item \textbf{数据点} $x_0$：从真实数据分布 $q(x_0)$ 中采样的干净样本
    \item \textbf{最终潜变量} $x_T$：标准高斯分布 $\mathcal{N}(0, I)$ 的样本
    \item \textbf{中间潜变量} $x_1, \dots, x_{T-1}$：介于数据和纯噪声之间的中间状态
    \item \textbf{前向过程}（Forward Process）：从 $x_0$ 到 $x_T$，逐步注入噪声
    \item \textbf{反向过程}（Reverse Process）：从 $x_T$ 到 $x_0$，逐步去噪
\end{itemize}

\begin{knowledgebox}{与 VAE 的关键区别}
在 VAE 中，编码器（后验分布）和解码器（似然分布）都需要学习。而在扩散模型中，前向过程是\textbf{预定义的}（不需要学习），我们只需要学习反向过程。此外，所有潜变量 $x_t$ 与数据 $x_0$ 具有\textbf{相同的维度}，这与 VAE 中潜变量通常维度更低不同。
\end{knowledgebox}

关于马尔可夫性质：前向过程和反向过程都被定义为马尔可夫过程，即 $x_t$ 仅依赖于 $x_{t-1}$（前向）或 $x_{t+1}$（反向），不依赖于更早的状态。形式上：

$$q(x_t | x_{t-1}, x_0) = q(x_t | x_{t-1})$$

这一性质大大简化了后续的数学推导。

\subsection{本章小结}

扩散模型通过预定义的前向噪声注入过程和可学习的反向去噪过程，在保持训练简单性的同时获得了强大的表达能力。与 VAE 不同，扩散模型不需要学习编码器，且所有潜变量与数据同维。

%% ========== 第二节：前向过程 ==========

